\scanpage{77}
Or les anneaux locaux des pts de $U$ sont
de dim \uline{$i\leq n$} ! Donc de $=$ les anneaux
strictement locaux $B'$ de $U$, et \deleted{\unsure{de}} $V$ étant
affine de dim. relative 1 sur $U$, donc
$V'=\unsure{V|B}\otimes_B B'$ \inserted{affine et} de dim relative 1 sur $B'$, on
\unsure{conçoit|constate|consort} que l'hyp $1'_n$ va impliquer
que $H^{\deleted{\unsure{p}}q}(V',F')=0$ si $q>i+1$ i.e. \uline{$i<q-1$},
[i.e. $R^qf_*(F)$ est nul en codim $i<q-1$].
C'est ce qu'il faut vérifier, \unsure{pièce|prière|rien} \unsure{en|au}

\uline{Lemme 3} $1'_n$ \deleted{\unsure{étant}} \inserted{[\deleted{\unsure{est}} pour des ess. t.f. sur corps]}
implique que
pour tt préschéma affine $X$, \unsure{ess.t} de
t. f. sur un corps sép. clos $k$, $X$ de
dim de Krull $d\leq n+1$, on a $\operatorname{cd}_{\ell}(X)\leq\dim X$

\uline{Dém} Par passage à la limite, on peut supposer
que $X$ est de t. f. sur $k$. \inserted{Récurrence sur $d=\dim X$, $d\leqn+1$} \deleted{Par lemme}
\inserted{Trivial si $d\leq0$. \deleted{Par lemme} soit donc $d\geq1$.} Par lemme
de normalisation, on peut supposer $X={\unsure{\mathbb E|\mathrm E}}^d_k$.
\inserted{($d\leq n+1$) : prouver $\operatorname{cd}_{d}(\mathbb E^d_k)\leq d$}
\deleted{Récurrence sur $d\leq n+1$.} Trivial si $d\leq0$.

En \unsure{projetant|général}, $X=\mathbb E^d_k\mathrel{\overset{f}{\unsure{-|\longrightarrow}}}Y=\mathbb E^{d'\deleted{\unsure{-1}}}$, où $d'=d-1\leq n$.
En vertu de $1'_n$, $R^qf_*(F)_{\bar y}\ne0$
$\Rightarrow\deleted{\unsure{dim}}\ q\leq\dim\underline O_{Y,\bar y}+1\deleted{\unsure{=q+2|i+1}}\inserted{=i+1}$. i.e. $R^q f_*(F)_{\bar y}$ est
% archive sheet number: 76
